\section{从插值多项式到梯形与 Simpson 求积}
\label{sec:08-Numerical-Analysis/04-Numerical-Differentiation-and-Integration/02}

你需要算 \(\int_0^1 e^{-x^2}\,dx\)。不定积分没有初等形式——大一微积分课上老师就告诉过你"这个积不出来"。但数值上你只需要一个六位有效数字的答案，不需要原函数。

你面前只有函数在离散点上的值。你觉得应该用梯形法——把区间切成很多小段，每段上把曲线近似成直线，面积就是梯形。公式你在别处见过：\(h[f_0/2+f_1+\cdots+f_{n-1}+f_n/2]\)。但你又隐约记得还有个 Simpson 法，据说是把曲线近似成抛物线，精度高出不少。这两个公式看着完全不像——梯形权重是 \(1/2,1,1,\ldots,1,1/2\)，Simpson 权重是 \(1,4,2,4,2,\ldots,4,2,4,1\)——它们之间有什么联系吗？

本节要建立一个统一的视角：梯形法和 Simpson 法都在做同一件朴素的事——先用容易积分的低次多项式替代原函数，再把替代函数精确积掉。从这个视角看去，两者的权重来源、误差阶数、乃至更高阶方法的推广路径，全部清晰可见。

\subsection{求积公式的通用模板：插值，再积分}

要近似

\[
I=\int_a^bf(x)\,dx,
\]

在 \([a,b]\) 上选 \(n+1\) 个节点 \(x_0,\ldots,x_n\)，做插值多项式 \(p_n(x)=\sum_{i=0}^nf(x_i)\ell_i(x)\)。因为插值多项式容易积分——它就是多项式的线性组合——直接对两边积分：

\[
\int_a^bf(x)\,dx\approx\int_a^bp_n(x)\,dx
=\sum_{i=0}^nf(x_i)\underbrace{\int_a^b\ell_i(x)\,dx}_{w_i}.
\]

每个节点 \(x_i\) 获得一个权重 \(w_i\)，这个权重就是把对应的 Lagrange 基函数从 \(a\) 到 \(b\) 积分得到的一个数。权重只依赖节点位置和积分区间，与函数值 \(f(x_i)\) 无关。

这个模板的统一性值得强调：你不需要单独记忆梯形法是一种面积近似、Simpson 法是另一种——两者都是"选节点 \(\to\) 插值 \(\to\) 积分基函数 \(\to\) 得到权重"这个流水线的输出。不同之处仅在于用了几个节点、放在什么位置。把求积公式放回"先插值再积分"这个框架后，所有公式的系数都不再是来路不明的魔法数字。

\subsection{梯形法：用弦替代曲线，积分弦}

单区间 \([a,b]\) 的情形最简单。两个端点的线性插值就是连接 \((a,f(a))\) 和 \((b,f(b))\) 的弦。Lagrange 基函数：

\[
\ell_0(x)=\frac{b-x}{b-a},\qquad
\ell_1(x)=\frac{x-a}{b-a}.
\]

分别积分：

\[
w_0=\int_a^b\frac{b-x}{b-a}\,dx=\frac{b-a}{2},\qquad
w_1=\int_a^b\frac{x-a}{b-a}\,dx=\frac{b-a}{2}.
\]

于是

\[
T=\frac{b-a}{2}[f(a)+f(b)].
\]

就是梯形的面积公式——弦下方直角梯形的面积。

将 \([a,b]\) 均匀分成 \(n\) 段，步长 \(h=(b-a)/n\)，每个子区间独立使用梯形公式，相邻区间在共享端点处的权重相加。内点被两个相邻区间各贡献一次半权重权，首次出现时"借走"半份，下一次出现时"归还"半份，合起来正好是一份完整权重；首尾端点只属于一个区间，拿半份。复合梯形法：

\[
T_n=h\left[
\frac12f(x_0)+\sum_{i=1}^{n-1}f(x_i)+\frac12f(x_n)
\right].
\]

若 \(f\in C^2[a,b]\)，由每个子区间上线性插值误差的积分，存在 \(\xi\) 使

\[
I-T_n=-\frac{b-a}{12}h^2f''(\xi).
\]

全局二阶：网格点数加倍，误差约四分之一。对凸函数，每个子区间上弦都在函数图像上方（凸函数的图像弯在弦之下），梯形法高估。例如 \(f(x)=x^2\) 在 \([0,1]\) 上单梯形给出 \(\frac12\)，真值 \(\frac13\)。

非均匀节点不能套用统一的 \(h\) 权重，应逐区间用各自的宽度；对光滑周期函数且端点各阶导数匹配，复合梯形法的误差中所有带端点导数的项会反复抵消，实际精度可能远高于 \(O(h^2)\)——这个"梯形法的超收敛"只对周期光滑函数成立，不是普遍性质。

\subsection{Simpson 法：两个小区间共享一条抛物线，对称性白送一阶}

现在把眼光放宽到三个等距点：\(a\)、中点 \(m=(a+b)/2\)、\(b\)。在这三点上做二次 Lagrange 插值，积掉三条基函数：

\[
\ell_0(x)=\frac{(x-m)(x-b)}{(a-m)(a-b)},\quad
\ell_1(x)=\frac{(x-a)(x-b)}{(m-a)(m-b)},\quad
\ell_2(x)=\frac{(x-a)(x-m)}{(b-a)(b-m)}.
\]

换元 \(x=a+th\)，在 \(t\in[0,2]\) 上计算积分，得到权重：

\[
w_0=\frac{h}{3},\qquad w_1=\frac{4h}{3},\qquad w_2=\frac{h}{3}.
\]

因为 \(b-a=2h\)，公式常写为

\[
S=\frac{b-a}{6}
[f(a)+4f(m)+f(b)].
\]

中点权重是中点的 4 倍——它同时服务于左右两半段，抛物线的顶点信息被加倍征用。

均匀网格上复合 Simpson 法要求子区间数 \(n\) 为偶数（每两个相邻子区间共享一组三个节点做一条抛物线）：

\[
S_n=
\frac h3\left[
f_0+f_n
+4\sum_{i\text{ odd}}f_i
+2\sum_{i\text{ even},\,0<i<n}f_i
\right].
\]

奇下标权重 \(4\)，内部偶下标权重 \(2\)——每对相邻子区间中，中点属于本抛物线的内部节点，独占对应基函数的全部积分；共享端点则是左右两条抛物线各贡献一半，合起来权重 \(2\)。

虽然由二次插值推导，Simpson 规则因对称性能精确积分三次多项式。理由如下：奇函数（如 \((x-m)^3\)）在 \([a,b]\) 上关于中点 \(m\) 积分为零；Simpson 的三点对称放置让三次奇函数对应的插值多项式精确捕捉到这一对称性。误差中 \(f'''\) 的贡献因此消失，第一个残留项来自 \(f^{(4)}\)。若 \(f^{(4)}\) 连续，

\[
I-S_n=-\frac{b-a}{180}h^4f^{(4)}(\xi),
\]

所以网格点数加倍，渐近误差约缩小 \(1/16\)。这和梯形法每加倍只缩小到 \(1/4\) 相比，是质的飞跃——同样的计算代价买到两个额外精度阶。

不规则网格上 \(1,4,2,\ldots,4,1\) 的权重模式不能照搬——等距是抛物线积分系数的前提。函数有尖点、跳跃或端点奇性时，高阶导数的存在性被破坏，四阶的误差结论退化回更低的实际阶数。

\subsection{Richardson 外推：梯形和 Simpson 本是一条渐近链上的相邻节点}

考虑梯形法用步长 \(h\) 和半步长 \(h/2\) 计算两次。如果梯形误差有渐近展开

\[
T_h=I+Ch^2+O(h^4),
\]

那么半步长为

\[
T_{h/2}=I+C\frac{h^2}{4}+O(h^4).
\]

现在你手里有两个近似值，已知其主导误差项之间恰有 4 倍关系。把第二个式子乘以 4，减第一个：

\[
\frac{4T_{h/2}-T_h}{3}
=I+\frac{4O(h^4)-O(h^4)}{3}
=I+O(h^4).
\]

\(h^2\) 阶的主误差项被消灭，剩余误差从二阶直接跳到四阶。这个组合恰好就是复合 Simpson 公式——不是近似意义上的"差不多"，而是代数上严格相等。

验证一下权重算术。记细网格步长 \(h'=h/2\)。粗网梯形 \(T_h\) 涉及点 \(x_0,x_2,x_4,\ldots\)（每隔一个点），细网梯形 \(T_{h/2}\) 涉及全部点 \(x_0,x_1,x_2,\ldots\)。在 \(\frac{4T_{h/2}-T_h}{3}\) 中：

\begin{itemize}
\tightlist
\item
  端点（如 \(x_0\)）：细网梯形贡献 \(4\cdot\frac{h'}{2}=2h'\)，粗网梯形贡献 \(-\frac{1}{3}\cdot h=-\frac{2h'}{3}\)，合为 \(\frac{4h'}{3}=\frac{h}{3}\cdot\frac{1}{2}\)，对应 Simpson 的端点权重 1；
\item
  奇下标内部点（只出现在细网）：\(4\cdot h'=\frac{4h}{3}\cdot\frac{1}{2}\)，对应 Simpson 权重 4；
\item
  偶下标内部点（两网共享）：细网 \(4\cdot\frac{h'}{2}=2h'\)，粗网 \(-\frac{1}{3}\cdot h=-\frac{2h'}{3}\)，合为 \(\frac{4h'}{3}\)，再除以 \(\frac{h}{3}\) 对应 Simpson 权重 2。
\end{itemize}

全部权重与前面推导的 \(S_n\) 完全一致。Richardson 外推之所以有效，是因为它利用了你对误差幂次的知识——已知主导误差正比于 \(h^2\)，就用两个分辨率的组合将其消去。同样的逻辑可以继续向上：消去 \(h^4\) 项得到更高阶的求积公式（Romberg 积分），每一步都在消耗一个渐近幂次。

用 \(n\) 和 \(2n\) 的结果比较来估计误差是同样的原理，但前提是已经进入预期的渐近区。如果两次近似共同遗漏一个窄峰——比如函数在 \(x=0.713\) 处有一个宽度远小于 \(h\) 的尖峰——\(T_n\) 和 \(T_{2n}\) 可能数值非常接近，却同时错得离谱。外推帮不了你看见从未被采样到的结构。

\subsection{积分比微分温和，但温和有温和的陷阱}

对比上一节数值微分的 U 形误差曲线，积分的处境要友善得多。局部函数值误差在积分中乘上区间宽度后求和——宽度 \(h\) 本身就是一个小数，把误差压缩了一阶。微分则相反：相减后还要除以小 \(h\)，把误差膨胀一阶。从信号处理的角度，积分是平滑操作（把高频幅度除以频率），微分是放大高频的操作（把高频幅度乘以频率）。

因此"积分比微分温和"是一个可以定量信赖的陈述。但温和不意味着万能：

\begin{itemize}
\tightlist
\item
  端点奇性。如果被积函数在积分限处发散——例如 \(\int_0^1 x^{-1/2}dx\) 在 \(x=0\) 处——等距节点的多项式插值在奇点附近完全失效，需要专门处理（如变量替换、奇异减法、或 Gauss--Jacobi 求积）。
\item
  强振荡。\(\int_0^{2\pi}\sin(100x)f(x)dx\) 型的积分，正负值快速交替，等距采样可能需要远超直觉的节点数才能分辨。Filon 方法和 Levin 型方法为这类问题专门设计。
\item
  极窄峰。带宽远小于网格间距的尖峰可能完全漏过所有采样点，得到精确而完全错误的零。自适应求积通过局部误差估计检测到"这里还需要更细"。
\item
  无限区间。\(\int_{-\infty}^\infty\) 需要截断、变量压缩（如 \(x=\tan t\)）或 Gauss--Hermite 型求积。
\item
  高维积分。一维的 \(n\) 点求积公式推广到 \(d\) 维变成 \(n^d\) 点——维数灾难直接作用在求积节点数上。稀疏网格、Monte Carlo 和拟 Monte Carlo 方法为高维而存在。
\end{itemize}

梯形法和 Simpson 法是你工具箱里最常被调用的两把通用工具。记住它们的来源——插值再积分——你就能随时推导出非均匀网格上的对应公式、判断何时精度阶会退化、以及决定是否值得向 Romberg 外推或自适应求积迈出下一步。
